% ============================================================
%  Abstract (ภาษาอังกฤษ + ภาษาไทย)
% ============================================================

% ── English Abstract ─────────────────────────────────────────
\clearpage
\phantomsection
\addcontentsline{toc}{chapter}{Abstract}

\begin{tabular}{@{}lp{10.5cm}@{}}
  Project Title   & An Exploration of Artificial Intelligence for Detection of \\
                  & Inhibition Zones in Antibiotic Susceptibility Testing \\[0.3cm]
  Credits         & 3 \\[0.2cm]
  Member(s)       & Mr. Pheeraphat Siriphatthanakhun \\
                  & Mr. Phattaradon Thanaprasitpattana \\[0.3cm]
  Project Advisor & Asst.\ Prof.\ Dr.\ Nontaphat Leerach \\[0.2cm]
  Co-advisor      & Asst.\ Prof.\ Dr.\ Paniti Achararit \\
                  & Dr.\ Karittha Jangsamsi \\[0.2cm]
  Supervisor (WIL)& Dr.\ Jirapan Premsuriya \\
                  & Dr.\ Jumphol Polvichai \\[0.3cm]
  Program         & Health Data Science \\[0.1cm]
  Department      & Computer Engineering \\[0.1cm]
  Faculty         & Engineering \\[0.1cm]
  Academic Year   & 2025 \\
\end{tabular}

\vspace{1cm}
\begin{center}\textbf{Abstract}\end{center}

\hspace{\parindent}%
Antibiotic resistance is a critical global health challenge that demands accurate diagnostic methods.
Manual measurement of zone of inhibition (ZOI) diameters in Kirby-Bauer disk diffusion testing---the internationally accepted Antimicrobial Susceptibility Testing (AST) standard---is time-consuming and prone to human error.
This project develops an AI-based web application to automate ZOI detection and measurement from photographs of antibiotic culture plates.

The system evaluates four ZOI detection approaches: a Convolutional Neural Network (CNN) baseline,
Faster R-CNN with a ResNet-50 backbone, YOLOv11n, and the classical Hough Circle Transform.
A Hybrid OCR pipeline---using Gemini 3.1 Flash lite as the primary engine and EasyOCR with preprocessing
(Otsu thresholding, adaptive binarization, and 24-angle rotation search) as a fallback---identifies
antibiotic disk labels. Detected zone diameters are interpreted against CLSI 2025 and EUCAST 2023 breakpoints
to classify bacteria as Susceptible (S), Intermediate (I), or Resistant (R).

A dataset of 287 real laboratory images was augmented to produce a training set of 2,442 images,
supplemented by 2,000 synthetic images split across train/validation/test partitions,
yielding a combined test-and-validation set of 235 images.
Experimental results indicate that deep learning-based approaches achieved better performance than the CNN classification baseline and classical Hough Circle Transform.
YOLOv11n achieved the best accuracy with an MAE of approximately 2.5--2.7\,mm and mAP@0.5 of 95.91\%,
outperforming Faster R-CNN (MAE 4.65\,mm), CNN (mAP@0.5 0.021), and Hough Transform (MAE $>$ 5\,mm).
Gemini 3.1 Flash lite achieved 71.72\% OCR accuracy on a comprehensive 198-label test set, outperforming Qwen2.5-VL:7b (49\%) and EasyOCR (28--32\%).
The complete system, including the React frontend, FastAPI backend, and SQLite database, was integrated into a deployable web application and evaluated in a real laboratory environment at the Princess Srisavangavadhana School of Medicine, Chulabhorn Royal Academy hospital server.
User evaluation by five domain experts yielded an overall usefulness score of 3.60 out of 5.00,
confirming the system's readiness for real-world microbiological laboratory workflows.

\noindent\textbf{Keywords:} Antimicrobial Susceptibility Testing, Zone of Inhibition, Deep Learning, Object Detection,
CNN, Faster R-CNN, YOLOv11, Hough Transform, Gemini, EasyOCR, Image Processing, CLSI, EUCAST, Web Application

\newpage

% ── บทคัดย่อ ─────────────────────────────────────────────────
\clearpage
\phantomsection
\addcontentsline{toc}{chapter}{บทคัดย่อ}

\begin{tabular}{@{}lp{10cm}@{}}
  หัวข้อปริญญานิพนธ์  & การสำรวจปัญญาประดิษฐ์สำหรับการตรวจจับโซนการยับยั้ง \\
                      & ในการทดสอบความไวต่อยาปฏิชีวนะ \\[0.3cm]
  หน่วยกิต            & 3 \\[0.2cm]
  ผู้เขียน             & นายพีรพัฒน์ ศิริพัฒนคุณ \\
                      & นายภัทรดนย์ ธนประสิทธิ์พัฒนา \\[0.3cm]
  อาจารย์ที่ปรึกษา     & ผศ.\ ดร.\ นนทพรรษน์ ลีราช \\[0.2cm]
  ที่ปรึกษาร่วม        & ผศ.\ ดร.\ ปณิธิ อัจฉราฤทธิ์ \\
                      & ดร.\ คริษฐา แจ้งสามสี \\[0.2cm]
  Supervisor (WIL)    & ดร.\ จิรพันธ์ เปรมสุริยา \\
                      & ดร.\ จุมพล พลวิชัย \\[0.3cm]
  หลักสูตร             & วิทยาศาสตรบัณฑิต \\[0.1cm]
  สาขาวิชา            & วิทยาศาสตร์ข้อมูลสุขภาพ \\[0.1cm]
  ภาควิชา             & วิศวกรรมคอมพิวเตอร์ \\[0.1cm]
  คณะ                 & วิศวกรรมศาสตร์ \\[0.1cm]
  ปีการศึกษา          & 2568 \\
\end{tabular}

\vspace{1cm}
\begin{center}\textbf{บทคัดย่อ}\end{center}

\hspace{\parindent}%
ปัญหาการดื้อยาต้านจุลชีพถือเป็นความท้าทายด้านสาธารณสุขระดับโลกที่ต้องการวิธีการวินิจฉัยที่แม่นยำ การวัดขนาดบริเวณยับยั้งเชื้อ (Zone of Inhibition: ZOI) ด้วยมือในการทดสอบ Kirby-Bauer disk diffusion---ซึ่งเป็นมาตรฐาน Antimicrobial Susceptibility Testing (AST) ที่ยอมรับในระดับสากล---ยังคงใช้เวลานานและมีโอกาสเกิดความผิดพลาดจากตัวบุคคล โครงงานนี้จึงพัฒนาเว็บแอปพลิเคชันที่ขับเคลื่อนด้วย AI เพื่อวิเคราะห์และวัดขนาด ZOI จากภาพถ่ายจานเพาะเชื้อแบบอัตโนมัติ

ระบบที่พัฒนาขึ้นประเมินแนวทางการตรวจจับ ZOI สี่วิธี ได้แก่ Convolutional Neural Network (CNN) แบบ Classification-based, Faster R-CNN (พร้อม ResNet-50 backbone), YOLOv11n และ Hough Circle Transform ในด้าน OCR ระบบใช้สถาปัตยกรรมแบบ Hybrid โดยมี Gemini 3.1 Flash lite เป็น Engine หลัก และ EasyOCR ร่วมกับ Otsu Thresholding และการค้นหามุมหมุน 24 มุม เป็น Fallback เพื่ออ่านรหัสยาปฏิชีวนะบนแผ่นยา ค่าเส้นผ่านศูนย์กลาง ZOI ที่ได้จะถูกนำไปเทียบกับเกณฑ์ CLSI 2025 และ EUCAST 2023 เพื่อจำแนกความไวต่อยาเป็น ไวต่อยา (S), กึ่งไว (I) หรือ ดื้อยา (R) โดยอัตโนมัติ

ชุดข้อมูลที่ใช้ประกอบด้วยภาพถ่ายจริง 287 ภาพ ผ่านกระบวนการ Data Augmentation จนได้ชุดฝึกสอนรวม 2,442 ภาพ และภาพสังเคราะห์อีก 2,000 ภาพที่แบ่งตามสัดส่วน 80/10/10 โดยชุดทดสอบและตรวจสอบรวมกันมี 235 ภาพ ผลการทดลองพบว่าโมเดล Deep Learning มีประสิทธิภาพเหนือกว่าทั้ง CNN และ Hough Transform อย่างชัดเจน โดย YOLOv11n ทำได้ดีที่สุดด้วยค่า MAE เพียง 2.5--2.7 มม. และ mAP@0.5 ถึง 95.91\% ทิ้งห่าง Faster R-CNN (MAE 4.65 มม.), CNN (mAP@0.5 0.021) และ Hough Transform (MAE $>$ 5 มม.) สำหรับ OCR พบว่า Gemini 3.1 Flash lite ให้ความแม่นยำสูงสุดที่ 71.72\% เหนือกว่า Qwen2.5-VL:7b (49\%) และ EasyOCR (28--32\%) ระบบทั้งหมดได้รับการบูรณาการเป็นเว็บแอปพลิเคชันที่พร้อมใช้งานจริง และได้รับการประเมินในสภาพแวดล้อมห้องปฏิบัติการจริงที่วิทยาลัยแพทยศาสตร์ศรีสวางควัฒน ราชวิทยาลัยจุฬาภรณ์ การประเมินความพึงพอใจจากผู้เชี่ยวชาญ 5 ท่านให้คะแนนรวม 3.60 จาก 5.00 ยืนยันความพร้อมของระบบสำหรับการใช้งานจริงในห้องปฏิบัติการจุลชีววิทยาคลินิก

\noindent\textbf{คำสำคัญ:} การทดสอบความไวต่อยาต้านจุลชีพ, บริเวณยับยั้งเชื้อ, Deep Learning, การตรวจจับวัตถุ,
CNN, Faster R-CNN, YOLOv11, Hough Transform, Gemini, EasyOCR, การประมวลผลภาพ, มาตรฐาน CLSI, มาตรฐาน EUCAST, เว็บแอปพลิเคชัน

\newpage
